% ══════════════════════════════════════════════════════════════════════════════
%  EXPERIMENT 2 — FRICTION AND MASS                    (Lesson 6, T3 Week 2)
%  From Lesson 6 of the A4 Forces Lesson Booklet; answer keys 6.* kept.
% ══════════════════════════════════════════════════════════════════════════════
\experimentdiv{Friction and Mass}{A fair test \textbullet{} Repeat trials and means \textbullet{} A line graph with a line of best fit}{6}\label{sec:2}
\lessoncodes{VC2S8U14 \textbullet{} VC2S8I01, I02, I03, I04, I05, I06}

\vocabhead
\begin{ctab}
\begin{tabular}{|SW{4.2cm}|W{11.6cm}|}
  \hline\hd & \thead{Word} & \thead{Definition (copy from the slide)} \\ \hline
  & \scaf{\textbf{Fair test}} & \scaf{An experiment where only one variable is changed and everything else is kept the same.} \\ \hline
  & \textbf{Independent variable} & \cell{6.v2} \\ \hline
  & \textbf{Dependent variable} & \cell{6.v3} \\ \hline
  & \textbf{Controlled variable} & \cell{6.v4} \\ \hline
  & \textbf{Trial} & \cell{6.v5} \\ \hline
  & \textbf{Mean} & \cell{6.v6} \\ \hline
\end{tabular}
\end{ctab}

\begin{lipanel}
\inquiryli{Plan a fair test, repeat each measurement and graph the means.}
\understandli{Find out how pressing two surfaces together harder changes friction.}
\sctitle
\begin{itemize}[leftmargin=1.4em, itemsep=0pt, topsep=1pt]
  \item I can name the variables and two controlled variables.
  \item I can calculate a mean from three trials.
  \item I can draw a line graph with a line of best fit and describe the trend with numbers.
\end{itemize}
\end{lipanel}

\begin{skillsfocus}
\focus{QP1}{Write a hypothesis in the faculty frame.}
\focus{PC2}{Name the variables, with two controlled.}
\focus{PC4}{Read the spring balance at a steady speed, three times.}
\focus{PA1}{Record and average the trials in a THUD table.}
\focus{PA2}{Plot the means and draw a line of best fit.}
\focus{PA3}{Calculate each mean.}
\focus{EV3}{Name an error and how to reduce it.}
\end{skillsfocus}

% ─── READING ─────────────────────────────────────────────────────────────────
\exppart{Before the lab}{Teal}

\begin{readingbox}{Reading --- Grip}
Friction is the force that makes walking possible. With every step, your shoe pushes backwards against the ground, and friction pushes you forwards. On ice, where friction is tiny, your foot slips backwards and you go nowhere.

Friction comes from the roughness of surfaces. Even surfaces that look smooth are covered in tiny bumps when you look at them under a microscope. When two surfaces press together, the bumps catch on each other and resist sliding. Press the surfaces together harder and the bumps push further into each other, so the friction is bigger.

Sometimes engineers want as much friction as possible --- in the rubber of a car tyre, the soles of running shoes or the grip tape on a skateboard. Other times they want as little as possible: oil in an engine and wax on skis reduce friction, so parts slide easily and less energy is wasted as heat.

In this experiment you will measure friction with a spring balance, then add masses to a wooden block so it presses harder on the bench. When you pull the block along at a \emph{steady} speed, your pull and the friction are balanced --- so the reading on the spring balance is equal to the friction force.
\end{readingbox}

\rcreset
\rcqk{6.c1}{Explain how friction helps you to walk, and why walking on ice is so hard.}{2.4cm}
\rcqk{6.c2}{Why does pressing two surfaces together harder make the friction bigger? Use the word \emph{bumps} in your answer.}{2.4cm}
\rcqk{6.c3}{Give one example where engineers want more friction, and one where they want less.}{2.2cm}
\rcqk{6.c4}{Why must the block move at a \emph{steady} speed when you read the spring balance?}{2.4cm}

% ─── TECHNIQUE ────────────────────────────────────────────────────────────────
\exppart{Technique spotlight}{Blue}

\begin{technique}{Making it a fair test}{PC2}
Change \textbf{one} thing (the \textbf{independent} variable), measure \textbf{one} thing (the \textbf{dependent} variable), and keep everything else the same (the \textbf{controlled} variables). Name at least two controlled variables.
\end{technique}

\begin{workedex}{Identifying variables}
A student rolls a toy car down a ramp at different heights and measures how far it travels across the floor.

\smallskip
\textbf{Independent:} the height of the ramp. \quad \textbf{Dependent:} the distance the car travels.\\
\textbf{Controlled:} the same car; the same floor surface; releasing the car without pushing it.
\end{workedex}

\begin{wedo}{Identifying variables}
A student times how long a spoonful of sugar takes to dissolve in water at 20\,\textdegree C, 40\,\textdegree C and 60\,\textdegree C.
\pt{a}Name the independent variable.
\abox{6.wd1a}{1.4cm}
\pt{b}Name the dependent variable.
\abox{6.wd1b}{1.4cm}
\pt{c}Name two controlled variables.
\abox{6.wd1c}{1.8cm}
\end{wedo}

\begin{workedex}{Calculating a mean}
With no mass added, three trials gave 0.8\,N, 1.0\,N and 0.9\,N.

\smallskip
Add the trials, then divide by how many there are: \quad $(0.8 + 1.0 + 0.9) \div 3 = 2.7 \div 3 = \mathbf{0.9\ \text{N}}$
\end{workedex}

\begin{wedo}{Calculating a mean}
With 400\,g on the block, a group measured 1.9\,N, 2.1\,N and 2.3\,N. Calculate the mean friction force.
\wbox{6.wd2}{2.4cm}
\end{wedo}

\begin{technique}{Join the points, or a line of best fit?}{PA2}
Every measurement has a little scatter, so measured results do not sit exactly on a line. Draw one smooth \textbf{line of best fit} through the middle of the points, with about as many points above it as below. Do \emph{not} join the points dot to dot: that line follows the scatter, not the trend.
\end{technique}

% ─── PLAN, SAFETY AND METHOD ─────────────────────────────────────────────────
\exppart{Plan, safety and method}{Teal}

\textbf{Aim:} To investigate the effect of changing the mass on a block on the friction force.

\task Complete the variables table. \skr{PC2}
\begin{ctab}
\begin{tabular}{|SW{4.4cm}|W{10.2cm}|}
  \hline\hd & \thead{Variable} & \thead{In this experiment} \\ \hline
  & \scaf{Independent} & \scaf{The mass added to the block (0 to 1000\,g)} \\ \hline
  & Dependent & \cell{6.t1-r2} \\ \hline
  & Controlled & \cell{6.t1-r3} \\ \hline
  & Controlled & \cell{6.t1-r4} \\ \hline
\end{tabular}
\end{ctab}

\needspace{8\baselineskip}
\task Write your hypothesis. \skr{QP1}

\scaf{If the mass on the block is increased, the friction force will \dots\ because \dots}
\abox{6.pred}{2.5cm}

\subsubsection*{Equipment}
Wooden block with a hook \textbullet{} spring balance (spring scale) or force sensor \textbullet{} five 200\,g masses

\begin{center}
\diagExpTwoSetup
\expfigcap{Pull the spring balance level with the bench, and read it while the block slides at a steady speed.}
\end{center}

\needspace{12\baselineskip}
\task Complete the risk assessment. \skr{PC3}
\begin{ctab}
\begin{tabular}{|SW{4.0cm}|W{5.2cm}|W{5.4cm}|}
  \hline\hd & \thead{Risk (what could go wrong)} & \thead{Prevention} & \thead{Responding (if it happens)} \\ \hline
  & \scaf{An overstretched spring balance snaps back} & \scaf{Pull gently and keep the balance in line with the block} & \scaf{Stop; check no one is hurt; tell the teacher} \\ \hline
  & Masses slide off the block and fall & & \\ \hline
\end{tabular}
\end{ctab}

\subsubsection*{Method}
\begin{enumerate}[itemsep=2pt]
  \item Place the wooden block on the bench top and attach the spring balance to its hook.
  \item Pull the block slowly, so it keeps moving at a steady speed. Read the force on the spring balance --- this is equal to the friction. Record it in Table 2.1.
  \item Repeat step 2 twice more, so you have three trials.
  \item Add a 200\,g mass on top of the block. Measure the friction three times and record the results.
  \item Repeat with 400, 600, 800 and 1000\,g on the block.
\end{enumerate}

% ─── RESULTS ─────────────────────────────────────────────────────────────────
\exppart{Results}{Amber}

\subsubsection*{Table 2.1: Friction force on the block for each added mass \skr{PA1}}
\begin{ctab}
\begin{tabular}{|SW{3.6cm}|Z{2.4cm}|Z{2.4cm}|Z{2.4cm}|Z{2.8cm}|}
  \hline\hd & \thead{Mass added (g)} & \thead{Friction, trial 1 (N)} & \thead{Friction, trial 2 (N)} & \thead{Friction, trial 3 (N)} & \thead{Mean (N)} \\ \hline
  & \scaf{e.g.\ 0 (block only)} & \scaf{0.8} & \scaf{1.0} & \scaf{0.9} & \scaf{0.9} \\ \hline
  & 0 (block only) & \cell{6.res-r2c2} & \cell{6.res-r2c3} & \cell{6.res-r2c4} & \cell{6.res-r2c5} \\ \hline
  & 200 & \cell{6.res-r3c2} & \cell{6.res-r3c3} & \cell{6.res-r3c4} & \cell{6.res-r3c5} \\ \hline
  & 400 & \cell{6.res-r4c2} & \cell{6.res-r4c3} & \cell{6.res-r4c4} & \cell{6.res-r4c5} \\ \hline
  & 600 & \cell{6.res-r5c2} & \cell{6.res-r5c3} & \cell{6.res-r5c4} & \cell{6.res-r5c5} \\ \hline
  & 800 & \cell{6.res-r6c2} & \cell{6.res-r6c3} & \cell{6.res-r6c4} & \cell{6.res-r6c5} \\ \hline
  & 1000 & \cell{6.res-r7c2} & \cell{6.res-r7c3} & \cell{6.res-r7c4} & \cell{6.res-r7c5} \\ \hline
\end{tabular}
\end{ctab}

\needspace{20\baselineskip}
\task Draw a line graph of your results. Put \textbf{mass added (g)} on the horizontal axis and \textbf{mean friction force (N)} on the vertical axis. Plot each point, draw a line of best fit, and give your graph a title. \skr{PA2}
\graphgrid{9cm}

% ─── ANALYSIS ────────────────────────────────────────────────────────────────
\exppart{Analysis and evaluation}{Coral}

\task Describe the trend in your results. Use the means for 0\,g and 1000\,g in your answer. \skr{PA5}
\abox{6.t2}{3.5cm}

\needspace{8\baselineskip}
\task Another group measured 2.1\,N, 2.3\,N and 2.2\,N with 600\,g on the block. Calculate the mean. \skr{PA3}
\wbox{6.t4}{2.6cm}

\needspace{8\baselineskip}
\task A group measured 1.6\,N, 1.7\,N and 2.6\,N with 400\,g on the block. Which value looks wrong? Find the mean of the other two.
\wbox{6.t5}{2.6cm}

\needspace{10\baselineskip}
\task Use the idea of \emph{bumps} pressing together to explain your results. Then describe a situation in everyday life where you have noticed this link between friction and mass.
\abox{6.t6}{4cm}

\needspace{14\baselineskip}
\task \skr{EV3}
\pt{a}Why did you take three trials for each mass?
\abox{6.t3a}{2.2cm}
\pt{b}Identify one error that could have affected your results.
\abox{6.t3b}{2.2cm}
\pt{c}Propose one improvement to the method that would reduce it.
\abox{6.t3c}{2.2cm}

\needspace{12\baselineskip}
\task Write a conclusion: state the trend, back it with numbers from your table, and say whether your hypothesis was supported.

\scaf{As the mass added increased from 0\,g to 1000\,g, the mean friction force increased from \dots\,N to \dots\,N. This supports / does not support the hypothesis.}
\abox{6.concl}{5cm}

\begin{yr8box}
\begin{keepext}
\exttask Look at your graph. With no masses added, the friction was not zero. Explain why, and say what your line of best fit would do if you could make the block lighter and lighter.
\abox{6.x8}{4cm}
\end{keepext}
\begin{keepext}
\exttask Predict how your results would change if the block slid on sandpaper instead of the bench. Would the line on your graph be steeper or less steep? Explain.
\abox{6.x9}{4cm}
\end{keepext}
\end{yr8box}

% ─── SKILL CHECK ─────────────────────────────────────────────────────────────
\exppart{Skill check}{NavyBlue}

\begin{skillcheck}
\scrow{QP1}{My hypothesis names the mass and the friction force.}
\scrow{PC2}{I named the variables, with two controlled.}
\scrow{PC4}{I read the balance at a steady speed, three times for each mass.}
\scrow{PA1}{My table has a title, headings, units in the headings and matching decimals.}
\scrow{PA2}{I drew a line of best fit, not dot to dot.}
\scrow{PA3}{I calculated each mean correctly.}
\scrow{EV3}{I named an error and a way to reduce it.}
\end{skillcheck}

\task Which skill will you focus on in Experiment 3? Explain what you will do differently.
\abox{6.sc}{5cm}
